\chapter*{编者说明}
\addcontentsline{toc}{chapter}{编者说明}
\markboth{编者说明}{编者说明}

\section*{一、本集的由来}

Strominger--Yau--Zaslow 猜想（下称 SYZ 猜想）自 1996 年提出以来，深刻改变了人们对镜像对称的理解：它把"镜像"从同调与量子上同调的对应，还原为一个可以几何化的图景------在 Calabi--Yau 流形的大复结构极限附近，存在由特殊拉格朗日子流形构成的环面纤维化，互为镜像的两族流形通过纤维上的 T 对偶相联系。此后二十余年，这一图景生长出特殊拉格朗日几何、仿射几何与热带几何、Gross--Siebert 对数退化纲领、非阿基米德几何与度量坍缩理论等多条研究脉络。

本集把六篇具有代表性的英文综述译为中文，合成一部体系化的译注本，目的是为中文读者提供一条从入门到前沿的连续阅读路径。

\section*{二、六篇综述与阅读建议}

\begin{enumerate}
  \item \textbf{Chan (2014)}，\emph{The Strominger--Yau--Zaslow Conjecture and Its Impact}：从 SYZ 猜想出发的完整脉络，篇幅短、框架清楚，是入门首选。
  \item \textbf{Chan (2016)}，\emph{A Glimpse of the SYZ Conjecture and Related Developments}（ICCM 2016 报告）：以"SYZ 之前的镜像对称"作铺垫，可与第一篇对照阅读。
  \item \textbf{Gross (2008)}，\emph{The SYZ Conjecture: From Torus Fibrations to Degenerations}：从环面纤维化讲到退化与对数结构，是 Gross--Siebert 纲领的起点。
  \item \textbf{Gross (2013)}，\emph{Mirror Symmetry and the SYZ Conjecture}：第三篇的全面升级版，把 2008--2013 年的进展纳入同一叙事；若只读一篇大综述，读这篇。
  \item \textbf{Yang Li (2022)}，\emph{Survey on the Metric SYZ Conjecture and Non-Archimedean Geometry}：度量侧的前沿综述，核心工具是多复变势理论与非阿基米德几何，并附公开问题。
  \item \textbf{Collins--Lin (2022)}，\emph{Recent Progress on SYZ Mirror Symmetry for Some Non-Compact Calabi--Yau Surfaces}：非紧 Calabi--Yau 曲面（有理椭圆曲面、del Pezzo 曲面）上的专题综述及其对引力瞬子分类的应用。
\end{enumerate}

建议的阅读顺序与本集章序一致；若时间有限，可先读第 1 章建立图景，再读第 4 章掌握系统叙述，第 5、6 章按研究兴趣选读。

\section*{三、译注本体例}

\begin{enumerate}
  \item \textbf{页码映射}：译文严格按原文页码推进，每页起始处在 LaTeX 源码中以 \texttt{\% [原文 p.N]} 标注（PDF 页号 = 印刷页号，偏移 0）。本集不改变原文的段落切分与论证顺序。
  \item \textbf{编号体系}：每篇综述构成一章（章号 = 论文序，1--6）；节号由 \LaTeX{} 自动编号为「章.节」，与原文节号一致；编号公式统一写作「章.原编号」（如第 1 章原式 (6.1) 记作 (1.6.1)），定理类对象记作「章.节.原编号」（如定理 4.5.2），图记作「章.图号」（如图 4.3）。正文中的交叉引用一律用纯文本。
  \item \textbf{批注}：译者的补充说明以 \texttt{\% [Tutor Note: ...]} 形式写在源码注释中，不进入排版；全部批注汇总于交付包中的《译注汇总》。批注只写可验证的易错符号、被省略的论证步骤、术语辨析与原文笔误，不添加与原文无关的内容。
  \item \textbf{原文笔误}：正文一律照原样排印；确认的笔误在批注中指出并给出正确形式，清单见交付包 README。
  \item \textbf{参考文献}：按学术惯例不翻译，保留英文原貌，由脚本从底本文字层自动转录（已知存在排印近似，见 README「已知限制」）。
  \item \textbf{术语}：全书使用统一术语表（随源码交付的 \texttt{glossary.tsv}）；人名、模型名与通行缩写（SYZ、CFT、K3、ALH 等）保留原文写法。
\end{enumerate}

\section*{四、编译与源码}

本集用 Lua\LaTeX{}（\texttt{ctexbook} 文档类）编译三次生成 PDF；交付包内附 \texttt{src/}（全部 \LaTeX{} 源码与图片）、\texttt{build.ps1}（一键编译脚本）与真实编译日志 \texttt{compile.log}。源码可在任何装有 TeX Live 2026 与中文字体的机器上独立重编译，不依赖底本 PDF 以外的中间素材。

\section*{五、版权与免责}

六篇底本的版权归原作者与出版方所有；本译本为学习与研究用途的非商业性译注，不替代原文。译文与批注中的疏漏、错误由译者负责，欢迎读者指正。原文的学术判断与结论以原文为准；译者批注仅代表译者理解。
